\begin{corrige}{4-13}

On remarque déjà que l'on doit avoir, pour $x=0$, $S(0)=1$ (fonction de survie égale à 1 au temps 0).

\begin{enumerate}

\item Pour $h(x)=\lambda$, on a:
\begin{eqnarray*}
( \ln S(x) )' &=& -\lambda,\\
\ln S(x) &=& k_1 - \lambda x,\\
S(x) & = & c_1 \exp( -\lambda x ), \quad (c_1 = e^{k_1}),\\
S(x) & = & \exp( -\lambda x ), \quad (c_1 = S(0) = 1),\\
F_X(x) & = & 1 - \exp (-\lambda x ),\\
f(x) & = & \lambda \exp (-\lambda x ).
\end{eqnarray*}

On reconnait ici la loi exponentielle de paramètre $\lambda$.

\item Pour $h(x)=\lambda k (\lambda x)^{k-1}$, on a:
\begin{eqnarray*}
( \ln S(x) )' &=& -\lambda k (\lambda x)^{k-1},\\
\ln S(x) &=& k_2 - (\lambda x)^k,\\
S(x) & = & c_2 \exp( - (\lambda x)^k ), \quad (c_2 = e^{k_2}),\\
S(x) & = & \exp( -(\lambda x)^k ), \quad (c_2 = S(0) = 1),\\
F_X(x) & = & 1 - \exp (-(\lambda x)^k ),\\
f(x) & = & \lambda k (\lambda x)^{k-1}\exp (-(\lambda x)^k ).
\end{eqnarray*}

C'est une loi dite de Weibull, qui est très utilisée pour des modèles de survie.

\end{enumerate}

\end{corrige}
